\section{Related Work and Design Gap}
\label{sec:related}

\subsection{Capabilities and Delegated Authorization}

Capability systems bind authority to explicit objects and support delegation without relying on ambient principal privileges. Dodson's cyber-physical capability architecture and CHERI Macaroons combine architectural capabilities with cryptographic network tokens, allowing scoped authority to be constructed, delegated, and exercised in robotics and industrial-control demonstrators~\cite{dodson2021,dodson2020}. SINT extends this direction to agent-facing gateways through signed tokens, physical maxima, attenuation-only delegation, revocation, approval tiers, and audit logging~\cite{sint2026}. ZTPM represents physical impact and live context as typed runtime policy dimensions for agents that control robots~\cite{ranathunga2026}, whereas DARC expresses multi-agent handoffs as contracts with attenuation, obligations, lineage, and revocation closure~\cite{bamil2026}. These systems establish practical and formal foundations for delegated physical authority. They do not establish when field-wise attenuation remains a refinement after a runtime layer rewrites or defers commands and a plant evolves.

\subsection{Physics-Aware Authorization and Runtime Enforcement}

Physics-aware authorization and runtime enforcement constrain commands through different mechanisms. CPAC derives access-control decisions from information flow over a physical grid model~\cite{etigowni2016}, and TwinGridShield evaluates grid-agent actions against a deterministic network twin while exposing sensitivity to model mismatch~\cite{rafy2026}. Runtime enforcement instead mediates an event stream: security automata characterize suppressible policies~\cite{schneider2000}, edit automata add insertion and replacement~\cite{ligatti2005}, Smart I/O places mediation close to industrial I/O~\cite{pearce2020}, and Lanotte \emph{et al.} synthesize monitors that correct or suppress compromised-controller actions with soundness and progress guarantees~\cite{lanotte2023}. Baird \emph{et al.} address scalable composition of such policies~\cite{baird2024}. \zonecap{} does not replace these mechanisms or claim greater enforcement throughput. It specifies the delegated authorization relation and the output contract an enforcer must satisfy for the child's induced physical behavior to remain within the parent's.

\subsection{Industrial Zone/Conduit Modeling}

IEC~62443-3-2 uses zones and conduits as central abstractions in security-risk assessment for industrial automation and control systems~\cite{iec62443}. Recent work automates parts of that risk-assessment workflow from engineering and threat models~\cite{dasilva2025}, and an IEC~62443-oriented domain-specific language formalizes zone/conduit installations and validation rules~\cite{vaudey2024}. These contributions operate primarily at design time; they do not turn a directed zone/conduit grant into a delegatable runtime object with an associated physical refinement relation. \zonecap{} uses selected, published zone/conduit structures as frozen compiler inputs, augments them with an explicitly declared evaluation overlay, and connects the resulting capabilities to runtime emissions and plant reachability. Only the zone/conduit topology is source-derived; target assets are deterministic encodings, and the behavioral and contextual fields are a declared overlay. This use is neither an implementation of the complete standard nor evidence of IEC~62443 conformance or certification.

